% A1 Portrait Poster Template (beamerposter)
\documentclass[final]{beamer}

% Poster package
\usepackage[size=a1,orientation=portrait,scale=1.2]{beamerposter}

% Theme
\usetheme{default}

% Packages
\usepackage{graphicx}
\usepackage{qrcode}
\usepackage{textpos}
\usepackage{tikz}
\usetikzlibrary{positioning,calc}
\usepackage{booktabs}

% Lengths
\setlength{\TPHorizModule}{1cm}
\setlength{\TPVertModule}{1cm}
\setlength{\topmargin}{-2cm}
\setlength{\headsep}{0pt}
\setlength{\headheight}{0pt}

% Colors
\definecolor{darkblue}{RGB}{20,20,80}
\definecolor{lightgray}{RGB}{240,240,240}

\setbeamercolor{block title}{fg=white,bg=darkblue}
\setbeamercolor{block body}{fg=black,bg=lightgray}
\setbeamerfont{block title}{size=\Large}
\setbeamerfont{block body}{size=\large}
\setbeamerfont{title}{size=\Huge}

% Metadata
\title{Advanced Class Imbalance Handling: SMOTE vs. GANs vs. Augmentation}
\author{Hakan Körpe}
\institute{TU Dortmund -- AIML Course}

\begin{document}
\begin{frame}[t]{}
\framesetup{width=\paperwidth}

\vspace*{-1.0cm}

% =======================
% TOP TITLE BAR
% =======================
\begin{tikzpicture}[remember picture,overlay]
\node[anchor=north west, inner sep=0pt] at (current page.north west)
{
\begin{minipage}{\paperwidth}
\begin{beamercolorbox}[wd=\paperwidth,ht=5cm,dp=1cm,center]{block title}
{\fontsize{40}{48}\selectfont \textbf{Advanced Class Imbalance Handling: SMOTE vs. GANs}}\\[0.4em]
{\Large Hakan Körpe — TU Dortmund}
\end{beamercolorbox}
\end{minipage}
};
\end{tikzpicture}

% QR Code
\begin{tikzpicture}[remember picture,overlay]
\node[anchor=north east, xshift=-1.2cm, yshift=-1.2cm]
at (current page.north east)
{\color{white}\qrcode[height=5em]{https://github.com/hakankorpe/AIML-class-imbalance/tree/second-project}};
\end{tikzpicture}

\vspace*{3.5cm}

\begin{columns}[t,totalwidth=\textwidth]

% =============================================================================
% LEFT COLUMN
% =============================================================================
\begin{column}{0.48\textwidth}

\begin{block}{Abstract}
The handling of extreme class imbalance is examined using a binary MNIST task. Three data-level approaches are compared: SMOTE interpolation, geometric augmentation, and conditional GAN generation. Geometric augmentation improves feature learning but fails to produce usable minority detections under severe imbalance. SMOTE is best for minority classes with a simple geometric structure. For topologically complex digits, structural properties are better preserved by GAN-based generation.
\end{block}

\begin{block}{Problem Setup \& Methodology}
\textbf{Dataset:} A binary MNIST classification task (Digit 0 vs Rest) is considered, with a \textbf{1:50 class imbalance}.\\
\textbf{Minority Samples:} Approximately 120 images are available.\\
\textbf{Objective:} Minority detection quality is evaluated under extreme imbalance.

\vspace{0.5em}
\begin{enumerate}
\item \textbf{SMOTE:} Synthetic samples are generated via linear interpolation between minority neighbors.
\item \textbf{Geometric Augmentation:} Input images are rotated ($\pm15^\circ$) without increasing class count.
\item \textbf{Conditional GAN:} A deep generator is trained exclusively on minority-class samples.
\end{enumerate}
\end{block}

\begin{block}{Experiment 1: Baseline vs. SMOTE}
\centering
\begin{columns}
\begin{column}{0.6\textwidth}
\includegraphics[width=\linewidth]{images/mnist_smote_conf.png}
\end{column}
\begin{column}{0.35\textwidth}
\textbf{Result:} Near-perfect separation is achieved (\textbf{0.96 Recall}).\\[0.4em]
\textbf{Insight:} Due to the simple topology of the digit, linear interpolation is sufficient for generating decision-consistent samples.
\end{column}
\end{columns}
\end{block}

\begin{block}{Experiment 2: The Augmentation Paradox}
\centering
\begin{columns}
\begin{column}{0.6\textwidth}
\includegraphics[width=\linewidth]{images/mnist_aug_conf.png}
\end{column}
\begin{column}{0.35\textwidth}
\textbf{Result:} High ranking performance is observed, while \textbf{zero precision} is obtained.\\[0.4em]
\textbf{Explanation:} At the selected decision threshold, predictions are assigned exclusively to the majority class.
\end{column}
\end{columns}
\end{block}

\begin{block}{Experiment 3: GAN Generation}
\centering
\begin{columns}
\begin{column}{0.6\textwidth}
\includegraphics[width=\linewidth]{images/mnist_gan_conf.png}
\end{column}
\begin{column}{0.35\textwidth}
\textbf{Result:} Strong recall (0.88), noisier boundary.\\[0.4em]
\textbf{Challenge:} Learning full distributions from $\approx$100 samples.
\end{column}
\end{columns}
\end{block}

\end{column}

% =============================================================================
% RIGHT COLUMN
% =============================================================================
\begin{column}{0.48\textwidth}

\begin{block}{Why Precision Matters}
Under extreme imbalance, \textbf{accuracy and ROC-AUC are misleading}. Predicting the majority class yields high accuracy but \textbf{zero usable detections}. We therefore prioritize \textbf{PR-AUC}.

\vspace{0.4em}
\small \textit{All confusion matrices are shown at threshold 0.5, reflecting deployment-time decisions.}
\end{block}

\begin{block}{Visual Comparison: Simple vs. Complex Topology}
\centering
\includegraphics[width=\linewidth]{images/grand_comparison_0.png}

\small \textit{Digit 0:} SMOTE preserves topology; GAN adds variance.

\vspace{1em}

\includegraphics[width=\linewidth]{images/grand_comparison_8.png}

\small \textit{Digit 8:} GAN preserves crossings; SMOTE introduces artifacts.
\end{block}

\begin{block}{Quantitative Benchmarks (Decision Metrics)}
\centering
\vspace{0.4em}
\begin{tabular}{lccc}
\toprule
\textbf{Method} & \textbf{PR-AUC} & \textbf{ROC-AUC} & \textbf{Status} \\
\midrule
Baseline (Weights) & 0.41 & 0.9953 & Reference \\
\textbf{SMOTE} & \textbf{0.96} & \textbf{0.9997} & \textbf{Best} \\
Augmentation & 0.00 & 0.9979 & Fail \\
GAN (cGAN) & 0.88 & $\approx$0.99 & Competitive \\
\bottomrule
\end{tabular}

\vspace{0.4em}
\end{block}

\begin{block}{Conclusion}
\begin{enumerate}
\item SMOTE is shown to be effective for minority classes with simple geometric structure.
\item GAN-based generation is found to better preserve complex topological features.
\item Geometric augmentation alone is insufficient for correcting extreme class imbalance.
\end{enumerate}
\end{block}
\end{column}
\end{columns}

% =======================
% BOTTOM BAR
% =======================
\begin{tikzpicture}[remember picture,overlay]
\node[anchor=south west, inner sep=0pt] at (current page.south west)
{
\begin{minipage}{\paperwidth}
\begin{beamercolorbox}[wd=\paperwidth,ht=3.5cm,dp=0.5cm]{block title}
\end{beamercolorbox}
\end{minipage}
};

\node[anchor=south west, xshift=1.2cm, yshift=0.8cm]
at (current page.south west)
{\color{white}\includegraphics[height=2.5cm]{images/tu.png}};
\end{tikzpicture}

\end{frame}
\end{document}